
\documentclass{article}
\usepackage{amsfonts}

%%%%%%%%%%%%%%%%%%%%%%%%%%%%%%%%%%%%%%%%%%%%%%%%%%%%%%%%%%%%%%%%%%%%%%%%%%%%%%%%%%%%%%%%%%%%%%%%%%%%%%%%%%%%%%%%%%%%%%%%%%%%%%%%%%%%%%%%%%%%%%%%%%%%%%%%%%%%%%%%%%%%%%%%%%%%%%%%%%%%%%%%%%%%%%%%%%%%%%%%%%%%%%%%%%%%%%%%%%%%%%%%%%%
%TCIDATA{OutputFilter=LATEX.DLL}
%TCIDATA{Version=5.50.0.2960}
%TCIDATA{<META NAME="SaveForMode" CONTENT="1">}
%TCIDATA{BibliographyScheme=Manual}
%TCIDATA{Created=Tuesday, June 30, 2026 16:59:53}
%TCIDATA{LastRevised=Tuesday, June 30, 2026 17:12:40}
%TCIDATA{<META NAME="GraphicsSave" CONTENT="32">}
%TCIDATA{<META NAME="DocumentShell" CONTENT="Scientific Notebook\Blank Document">}
%TCIDATA{CSTFile=Math with theorems suppressed.cst}
%TCIDATA{PageSetup=72,72,72,72,0}
%TCIDATA{AllPages=
%F=36,\PARA{038<p type="texpara" tag="Body Text" >\hfill \thepage}
%}


\newtheorem{theorem}{Theorem}
\newtheorem{acknowledgement}[theorem]{Acknowledgement}
\newtheorem{algorithm}[theorem]{Algorithm}
\newtheorem{axiom}[theorem]{Axiom}
\newtheorem{case}[theorem]{Case}
\newtheorem{claim}[theorem]{Claim}
\newtheorem{conclusion}[theorem]{Conclusion}
\newtheorem{condition}[theorem]{Condition}
\newtheorem{conjecture}[theorem]{Conjecture}
\newtheorem{corollary}[theorem]{Corollary}
\newtheorem{criterion}[theorem]{Criterion}
\newtheorem{definition}[theorem]{Definition}
\newtheorem{example}[theorem]{Example}
\newtheorem{exercise}[theorem]{Exercise}
\newtheorem{lemma}[theorem]{Lemma}
\newtheorem{notation}[theorem]{Notation}
\newtheorem{problem}[theorem]{Problem}
\newtheorem{proposition}[theorem]{Proposition}
\newtheorem{remark}[theorem]{Remark}
\newtheorem{solution}[theorem]{Solution}
\newtheorem{summary}[theorem]{Summary}
\newenvironment{proof}[1][Proof]{\noindent\textbf{#1.} }{\ \rule{0.5em}{0.5em}}
\input{tcilatex}

\begin{document}


\begin{eqnarray*}
K_{rot} &=&\frac{1}{2}\mathbb{I\omega }^{2} \\
&=&\frac{1}{2}\left( 2\times 2.\,\allowbreak 693\,5\times 10^{-56}\unit{m}%
^{2}\unit{kg}\right) \left( 2000\frac{\unit{rad}}{\unit{s}}\right) ^{2} \\
&=&\allowbreak 1.\,\allowbreak 077\,4\times 10^{-49}\unit{J}
\end{eqnarray*}%
\[
K_{rot}=\frac{1}{2}k_{B}T=\frac{1}{2}\left( 1.3806568\times 10^{-23}\unit{J}%
\unit{K}^{-1}\right) \left( 293\unit{K}\right) 
\]%
$\frac{1}{2}\left( 1.3806568\times 10^{-23}\unit{J}\unit{K}^{-1}\right)
\left( 293\unit{K}\right) =\allowbreak 2.\,\allowbreak 022\,7\times 10^{-21}%
\unit{J}$%
\[
\frac{1}{2}\left( 2\times 2.\,\allowbreak 693\,5\times 10^{-56}\unit{m}^{2}%
\unit{kg}\right) \mathbb{\omega }^{2}=\allowbreak 2.\,\allowbreak
022\,7\times 10^{-21}\unit{J}
\]%
, Solution is: $\frac{2.\,\allowbreak 740\,4\times 10^{17}}{\unit{s}},-\frac{%
2.\,\allowbreak 740\,4\times 10^{17}}{\unit{s}}\allowbreak $,

\bigskip 

We should think about how big this energy would be for a single particle,
say, nitrogen atoms $\left( m=2.33\times 10^{-26}\unit{kg}\right) $. 
\[
K_{T}=\frac{1}{2}mv^{2}
\]%
and%
\[
v_{rms}=\sqrt{\frac{3\left( 1.381\times 10^{-23}\frac{\unit{J}}{\unit{K}}%
\right) \left( 293\unit{K}\right) }{\left( 2.33\times 10^{-26}\unit{kg}%
\right) }}=721.\,\allowbreak 79\frac{\unit{m}}{\unit{s}}
\]%
so%
\begin{eqnarray*}
K_{T} &=&\frac{1}{2}\left( 2.33\times 10^{-26}\unit{kg}\right) \left(
721.\,\allowbreak 79\frac{\unit{m}}{\unit{s}}\right) ^{2} \\
&=&\allowbreak 6.\,\allowbreak 069\,4\times 10^{-21}\unit{J}
\end{eqnarray*}%
which isn't very big. That is not a surprise. Atoms are small.

\end{document}
